\chapter{The Fourier seed}\label{ch:fourier-seed}


\index{Fourier seed|textbf}
\index{propagator!logarithmic|textbf}
\index{bar construction!Fourier seed}

The Fourier seed of the bar construction is the residue kernel
\[
\eta_{ij}=d\log(z_i-z_j)
\]
on the two-point collision divisor. It is not an analytic Fourier
transform on functions. It is the chain-level mechanism by which the
OPE of a chiral algebra~$\cA$ is converted into the differential of
the conilpotent bar coalgebra
$\barB_X(\cA)=T^c(s^{-1}\bar\cA)$, with logarithmic forms on the
Fulton--MacPherson compactifications~$\overline{C}_n(X)$ carrying the
configuration-space part.

Three genus-$0$ statements are exact:
\begin{enumerate}[label=\textup{(\roman*)}]
\item \emph{Bar coalgebra.}
 $\barB_X(\cA)_n=(s^{-1}\bar\cA)^{\boxtimes n}\otimes
 \Omega^\bullet(\overline{C}_n(X),\log D)$, with coproduct induced by
 deconcatenation and boundary factorization.
\item \emph{Propagator.} The bar differential $d = \sum_{\{i,j\}}
 \Res_{D_{ij}}(\eta_{ij}\otimes\mu_{ij})+d_{\mathrm{int}}+
 d_{\mathrm{dR}}$ extracts the OPE coefficient through the chiral
 product~$\mu$.
\item \emph{Kernel.} The identity $d^2 = 0$ is a formal consequence
 of the Arnold relation
 $\eta_{12}\wedge\eta_{23} + \eta_{23}\wedge\eta_{31} +
 \eta_{31}\wedge\eta_{12} = 0$
 in~$\Omega^2(\overline{C}_3(X), \log D)$.
\end{enumerate}
At genus~$g \geq 1$, the prime-form propagator has an Arakelov
correction. The cyclic Arnold form acquires a canonical cohomology
class; the algebra-dependent scalar multiplying the resulting genus
obstruction is the modular characteristic~$\kappa(\cA)$ on the
scalar lane of Theorem~\ref{thm:modular-characteristic}. Thus
$\kappa(\cA)$ measures the algebraic coupling to the geometric
curvature class, not the existence of the class itself.

% ================================================================
\section{The kernel}\label{sec:fourier-kernel}
% ================================================================

\begin{definition}[Logarithmic propagator]\label{def:fourier-log-propagator}
\index{propagator!logarithmic}
For distinct points $z_1, z_2$ on a smooth algebraic curve~$X$, the
\emph{logarithmic propagator} is
\begin{equation}\label{eq:fourier-log-propagator}
\eta_{ij} \;=\; d\log(z_i - z_j)
\;=\; \frac{dz_i - dz_j}{z_i - z_j}
\;\in\; \Omega^1(\overline{C}_2(X),\, \log D_{ij}).
\end{equation}
\end{definition}

\begin{proposition}[Three properties of the propagator;
 \ClaimStatusProvedHere]
\label{prop:fourier-propagator-properties}
The logarithmic propagator satisfies:
\begin{enumerate}
\item \emph{Simple pole.}
 $\eta_{ij}$ has a simple pole along
 $D_{ij} = \{z_i = z_j\}$ with
 $\Res_{D_{ij}}\eta_{ij} = 1$.
\item \emph{Closedness.}
 $d\eta_{ij} = 0$ on $C_2(X) \setminus D_{ij}$.
\item \emph{Arnold relation.}
 In $\Omega^2(\overline{C}_3(X),\, \log D)$,
 \begin{equation}\label{eq:fourier-arnold}
 \eta_{12} \wedge \eta_{23}
 + \eta_{23} \wedge \eta_{31}
 + \eta_{31} \wedge \eta_{12} = 0.
 \end{equation}
\end{enumerate}
\end{proposition}

\begin{proof}
Properties~(1) and~(2) are immediate: the pole is simple with
residue $\Res_{\epsilon=0}(d\epsilon/\epsilon) = 1$, and
$d\,d\!\log = 0$.

For~(3), set $u_{ij} = z_i - z_j$ and write
$\eta_{ij} = u_{ij}^{-1}\,du_{ij}$ where
$du_{ij} = dz_i - dz_j$.
Each wedge product
$\eta_{ij} \wedge \eta_{jk}
 = u_{ij}^{-1} u_{jk}^{-1}\, du_{ij} \wedge du_{jk}$.
Since $du_{12} + du_{23} + du_{31} = 0$, all three wedge products
share a common $2$-form
$N := du_{12} \wedge du_{23}
= du_{23} \wedge du_{31} = du_{31} \wedge du_{12}$.
(For instance:
$du_{23} \wedge du_{31}
= du_{23} \wedge (-du_{12} - du_{23})
= -du_{23} \wedge du_{12} = du_{12} \wedge du_{23} = N$.)
Then
\begin{align}
\sum_{\mathrm{cyc}} \eta_{12} \wedge \eta_{23}
&= \left[
 \frac{1}{u_{12}\,u_{23}}
 + \frac{1}{u_{23}\,u_{31}}
 + \frac{1}{u_{31}\,u_{12}}
\right] \cdot N
\notag \\
&= \frac{u_{31} + u_{12} + u_{23}}{u_{12}\,u_{23}\,u_{31}}
 \cdot N
= \frac{0}{u_{12}\,u_{23}\,u_{31}} \cdot N = 0.
\label{eq:fourier-arnold-computation}
\end{align}
The numerator vanishes because
$(z_3 - z_1) + (z_1 - z_2) + (z_2 - z_3) = 0$.
\end{proof}

\begin{proposition}[Genus-\texorpdfstring{$1$}{1} propagator;
\ClaimStatusConditional]
\label{prop:fourier-genus1-propagator}
On the elliptic curve
$E_\tau = \mathbb{C}/(\mathbb{Z} + \tau\mathbb{Z})$, the
propagator acquires a period correction:
\begin{equation}\label{eq:fourier-genus1-propagator}
\eta^{(1)}_{ij} = d\log\theta_1(z_i - z_j \,|\, \tau)
+ \frac{2\pi i}{\operatorname{Im}\tau}\,
 \operatorname{Im}(z_i - z_j)\,d\bar{z}_i,
\end{equation}
where $\theta_1(u|\tau)$ is the Jacobi theta function with simple
zero at $u = 0$. The correction descends the theta quotient to
$E_\tau \times E_\tau$. For
\[
\mathcal A^{(1)}_{123}:=
\eta^{(1)}_{12}\wedge\eta^{(1)}_{23}
+\eta^{(1)}_{23}\wedge\eta^{(1)}_{31}
+\eta^{(1)}_{31}\wedge\eta^{(1)}_{12},
\]
the canonical statement is the current/cohomology identity
\begin{equation}\label{eq:fourier-arnold-failure}
[\mathcal A^{(1)}_{123}]
=2\pi i\,[\omega_\tau],
\end{equation}
where
$\omega_\tau = \frac{i}{2\operatorname{Im}\tau}\,dz \wedge d\bar{z}$
is the normalized Arakelov $(1,1)$-form on $E_\tau$. A pointwise
equality of ordinary forms is not canonical without choosing the
Arakelov--Green representative and interpreting the diagonal
contribution as a logarithmic current.
\end{proposition}

\begin{proof}
The holomorphic part
$\partial_{z_i}\!\log\theta_1(z_i - z_j|\tau)\,dz_i$ has the
required simple pole at $z_i = z_j$ (since $\theta_1$ vanishes
simply at the origin) but picks up monodromy
$\partial_u\!\log\theta_1(u + \tau|\tau)
 = \partial_u\!\log\theta_1(u|\tau) - 2\pi i$
around the $B$-cycle. The correction term
$\frac{2\pi i}{\operatorname{Im}\tau}\operatorname{Im}(z_i - z_j)
\,d\bar{z}_i$
compensates this monodromy, making $\eta^{(1)}_{ij}$ single-valued.

Single-valuedness forces a non-holomorphic de~Rham component. In
Arakelov normalization the logarithmic current $d\eta^{(1)}_{ij}$ is
the difference between the diagonal residue current and the normalized
volume form; equivalently its off-diagonal cohomology class is
represented by $2\pi i\,\omega_\tau$ (Faltings~\cite{Fal84}).

The genus-$0$ Arnold cancellation used
$u_{12}+u_{23}+u_{31}=0$ inside a single affine coordinate. On
$E_\tau$ the same local cancellation survives near a triple
collision, while the period correction contributes the global class.
The normalized computation is precisely the elliptic clause of
Theorem~\ref{thm:arnold-higher-genus}:
\[
[\mathcal A^{(1)}_{123}]=2\pi i[\omega_\tau].
\]
The proposition records that theorem in the notation of this seed
chapter.
\end{proof}


% ================================================================
\section{The transform}\label{sec:fourier-transform}
% ================================================================

\begin{definition}[Bar complex at tensor power \texorpdfstring{$n$}{n}]
\label{def:fourier-bar-transform}
For a chiral algebra $\cA$ on a smooth curve $X$, the bar complex
at tensor power $n$ is
\begin{equation}\label{eq:fourier-bar-transform}
\barB_n(\cA) = (s^{-1}\bar{\cA})^{\otimes n}
\otimes \Omega^\bullet(\overline{C}_n(X),\, \log D)
\end{equation}
with differential $d = d_{\mathrm{int}} + d_{\mathrm{res}}
+ d_{\mathrm{dR}}$. The three components correspond to the internal
differential of $\cA$, residue extraction along collision divisors
$D_{ij} \subset \overline{C}_n(X)$, and the de~Rham differential on
logarithmic forms. The general construction is developed in
Chapter~\ref{chap:bar-cobar}; here we compute it.
\end{definition}

\subsection*{Computation~(a): $X = \operatorname{Spec} k$,
 $\cA = \Lambda(V)$}

\begin{computation}[Bar on a point]\label{comp:fourier-point}
Let $X = \operatorname{Spec} k$ (a point) and $\cA = \Lambda(V)$
(the exterior algebra on a finite-dimensional vector space $V$).
\begin{enumerate}
\item $\overline{C}_n(\operatorname{Spec} k) = \operatorname{Spec} k$
 for all $n$.
\item $\Omega^\bullet(\overline{C}_n,\, \log D) = k$: no forms, no
 propagator.
\item $d_{\mathrm{dR}} = 0$ and $d_{\mathrm{res}}$ reduces to the
 binary product differential.
\item $\barB(\Lambda(V)) = T^c(s^{-1}\bar{V},\, d_{\mathrm{res}})$,
 the classical bar construction of
 Loday--Vallette~\cite[\S2.2]{LV12}.
\end{enumerate}
\end{computation}

\begin{proposition}[\ClaimStatusProvedHere]
\label{prop:fourier-com-lie-duality}
For $\dim V = 1$, the bar complex has zero differential:
$\barB(\Lambda(V)) \cong \operatorname{Sym}^c(s^{-1}V)$. The
exterior algebra supplies the one-dimensional Koszul seed whose
finite dual is the symmetric coalgebra.
\end{proposition}

\begin{proof}
The reduced part $\bar{V} = V$ is one-dimensional, so
$(s^{-1}V)^{\otimes n} \cong k$ for all $n$. The bar differential
$d_{\mathrm{res}}\colon (s^{-1}v)^{\otimes 2} \mapsto
s^{-1}(v \wedge v) = 0$
vanishes because $v \wedge v = 0$ in $\Lambda(V)$. The cofree
coalgebra on a one-dimensional space is cocommutative, giving
$\operatorname{Sym}^c(s^{-1}V)$.
\end{proof}

\begin{remark}\label{rem:fourier-com-lie-operad}
At the operad level, $\Lambda(V)$ is a free $\mathsf{Com}$-algebra.
Its Koszul dual coalgebra is a cofree $\mathsf{Com}^!$-coalgebra, and
$\mathsf{Com}^! = \mathsf{Lie}$
(Loday--Vallette~\cite[Theorem~7.2.2]{LV12}). For $\dim V = 1$,
the Lie coalgebra structure is trivial, so the cofree Lie coalgebra
coincides with the symmetric coalgebra. For $\dim V > 1$, the
nontrivial Lie coalgebra differential encodes the shuffle relations.
\end{remark}

\subsection*{Computation~(b): $X = \mathbb{P}^1$,
 $\cA = \mathcal{H}_k$}

The Heisenberg chiral algebra $\mathcal{H}_k$ at level
$k \in \mathbb{C}^\times$ has a single current $\alpha(z)$ of
conformal weight~$1$ with OPE
\begin{equation}\label{eq:fourier-heisenberg-ope}
\alpha(z)\alpha(w) = \frac{k}{(z-w)^2} + \text{regular}.
\end{equation}
The double pole and the absence of a simple pole are the entire
structure.

\begin{computation}[\texorpdfstring{$n = 1$}{n = 1}]\label{comp:fourier-heisenberg-n1}
The bar complex at tensor power~$1$ is $\barB_1 = s^{-1}\alpha$.
There is no collision divisor, so
$d_{\mathrm{res}} = 0$ and $d_{\mathrm{dR}} = 0$. The element
$s^{-1}\alpha$ is a cocycle.
\end{computation}

\begin{computation}[\texorpdfstring{$n = 2$}{n = 2}; \ClaimStatusProvedHere]
\label{comp:fourier-heisenberg-n2}
At tensor power~$2$, the bar element is
\begin{equation}\label{eq:fourier-bar2-heisenberg}
\xi_{12} = s^{-1}\alpha(z_1) \otimes s^{-1}\alpha(z_2)
 \otimes \eta_{12}
\;\in\; \barB_2(\mathcal{H}_k).
\end{equation}
The bar differential acts by restricting to the boundary stratum
$D_{12} \subset \overline{C}_2(\mathbb{P}^1)$ and extracting the
OPE data via the Borcherds identity
(cf.~Chapter~\ref{ch:heisenberg-frame},
equation~\eqref{eq:frame-borcherds-extraction}).

\medskip\noindent\textbf{Step~1.}
The OPE $\alpha(z_1)\alpha(z_2) = k/(z_1 - z_2)^2 + \text{reg.}$
has a second-order pole with constant coefficient $c_2 = k$ and no
simple pole ($c_1 = 0$).

\medskip\noindent\textbf{Step~2.}
The boundary pairing on $D_{12}$ pairs the OPE singularity with the
logarithmic test form $\eta_{12}$. For a second-order pole with
coefficient $c_2$, the Borcherds extraction gives
\begin{equation}\label{eq:fourier-borcherds-extraction}
d_{\mathrm{res}}(\xi_{12})
= \bigl\langle c_2/(z_1 - z_2)^2,\; \eta_{12}
 \bigr\rangle_{D_{12}}
= c_2 = k.
\end{equation}

\medskip\noindent\textbf{Step~3.}
The result is $k \cdot 1 \in \barB_0 = \mathbb{C}$: the
second-order OPE coefficient is extracted into the degree-$0$
component of the bar complex.
\end{computation}

\begin{remark}[Borcherds mechanism]\label{rem:fourier-borcherds}
The extraction~\eqref{eq:fourier-borcherds-extraction} uses the
full chiral product~$\mu$
(Convention~\ref{conv:product-vs-bracket}),
which sums \emph{all} OPE modes (not just the simple pole).
This is \emph{not} the na\"ive contour integral
$\Res_{\epsilon = 0}[k\epsilon^{-2} \cdot d\epsilon/\epsilon]$
(which involves a triple pole and vanishes).
The chiral product extracts data via the Borcherds identity
(\S\ref{sec:frame-bar-deg1}), not by multiplicatively composing
the OPE pole with the logarithmic form.
\end{remark}

\begin{computation}[\texorpdfstring{$n = 3$}{n = 3}; \ClaimStatusProvedHere]
\label{comp:fourier-heisenberg-n3}
At tensor power~$3$, the bar element is
\[
\xi_{123} = s^{-1}\alpha(z_1) \otimes s^{-1}\alpha(z_2)
 \otimes s^{-1}\alpha(z_3)
 \otimes \eta_{12} \wedge \eta_{23}.
\]
Three boundary strata $D_{12}$, $D_{23}$, $D_{13}$ contribute to
$d_{\mathrm{res}}$, each extracting the second-order OPE coefficient:
\begin{align}
d_{\mathrm{res}}(\xi_{123})
&= k \cdot s^{-1}\alpha(z_3) \otimes
 \eta_{23}\big|_{z_1 = z_2}
\notag \\
&\quad + k \cdot s^{-1}\alpha(z_1) \otimes
 \eta_{12}\big|_{z_2 = z_3}
\notag \\
&\quad + k \cdot s^{-1}\alpha(z_2) \otimes
 \eta_{12}\big|_{z_1 = z_3}.
\label{eq:fourier-dres-n3}
\end{align}
Applying $d_{\mathrm{res}}$ again produces triple collisions.
The Arnold relation~\eqref{eq:fourier-arnold} forces the three-term
sum of triple residues to vanish:
$d_{\mathrm{res}}^2 = 0$ by the
cancellation~\eqref{eq:fourier-arnold-computation}.
\end{computation}

\begin{theorem}[Heisenberg bar seed; \ClaimStatusProvedHere]
\label{thm:fourier-heisenberg-bar}
\index{Heisenberg!bar complex}
The Koszul contraction of the bar complex of the Heisenberg algebra on
$\mathbb{P}^1$ is
\begin{equation}\label{eq:fourier-heisenberg-bar}
H^\bullet(\barB(\mathcal{H}_k)) \;\cong\;
\operatorname{Sym}^{\mathrm{ch},c}(V^*[-1])
\end{equation}
where $V = \langle\alpha\rangle$ is the one-dimensional space of
currents. The scalar $k$ is recorded in the bar differential by
$d_{\mathrm{res}}(\xi_{12})=k\cdot1$; after finite-type Verdier
duality it becomes the curvature $m_0=-k\,\omega$ of the Koszul-dual
commutative chiral algebra.
\end{theorem}

\begin{proof}
The bar complex has a single primitive class represented by
$s^{-1}\alpha$. The cofree coalgebra on one primitive is
cocommutative, giving the symmetric coalgebra in
\eqref{eq:fourier-heisenberg-bar}. The scalar is not lost: the
degree-$2$ extraction
$d_{\mathrm{res}}(\xi_{12}) = k \cdot 1$
(Computation~\ref{comp:fourier-heisenberg-n2}) is dual, after
Verdier finite duality, to the curvature term
$m_0=-k\,\omega$ of the commutative algebra
$\operatorname{Sym}^{\mathrm{ch}}(V^*)$.
The full verification is Theorem~\ref{thm:heisenberg-bar-complete}
in Chapter~\ref{ch:heisenberg-frame}.
\end{proof}

\begin{remark}[Chiral homology verification]
\label{rem:chiral-homology-heisenberg-verification}
\index{chiral homology!Heisenberg verification}
The Beilinson--Drinfeld chiral homology
$H^*_{\mathrm{ch}}(\mathbb{P}^1, \mathcal{H}_k)$ is, by
definition, the cohomology of the de~Rham complex of the
$\mathcal{D}$-module $\Delta_!(\mathcal{H}_k)$ on the Ran space
(\cite[\S4.1]{BD04}). For $\mathbb{P}^1$ with one current $J$ of
weight~$1$: $H^0_{\mathrm{ch}} = \mathbb{C}$ (vacuum coinvariant,
since $J_{-1}|0\rangle$ generates under the global residue).
Bar cohomology is $H^1(\bar{B}(\mathcal{H}_k)) = \mathbb{C}\cdot
\alpha^*$ (one-dimensional, by PBW concentration:
Theorem~\ref{thm:fourier-heisenberg-bar}) and
$H^n(\bar{B}(\mathcal{H}_k)) = 0$ for $n \geq 2$. This
confirms the BD identification $H^*_{\mathrm{ch}} = H^*(\bar{B})$
on the simplest non-trivial example: the bar complex is concentrated
in degree~$1$, reflecting chiral Koszulness.
\end{remark}

\begin{remark}[Heisenberg is not self-dual]
\label{rem:fourier-heisenberg-not-selfdual}
$\mathcal{H}_k^! =
(\operatorname{Sym}^{\mathrm{ch}}(V^*),m_0=-k\,\omega)$, not
$\mathcal{H}_{-k}$. Self-duality would require the Koszul dual to
be a Heisenberg algebra at the dual level. Instead,
$\mathcal{H}_k^!$ is the curved chiral symmetric algebra, a free
chiral commutative algebra with curvature, not a Heisenberg algebra
with a nontrivial OPE.
\end{remark}

\begin{proposition}[The five duality objects and the suspension;
\ClaimStatusProvedHere]
\label{prop:fourier-five-duality-objects}
Let $\cA$ be an augmented chiral algebra whose bar pieces are
finite-type holonomic $\mathcal D$-modules on $\Ran(X)$ after the
chosen Koszul contraction. Then the following objects have different
targets:
\begin{equation}\label{eq:fourier-five-objects}
\cA,\qquad
\barB_X(\cA)=T^c(s^{-1}\bar\cA),\qquad
\cA^i:=H^\bullet(\barB_X(\cA)),\qquad
\cA^!:=\mathbb{D}_{\Ran}(\cA^i),\qquad
Z^{\mathrm{der}}_{\mathrm{ch}}(\cA).
\end{equation}
The bar-cobar counit
\begin{equation}\label{eq:fourier-counit-not-dual}
\Omega_X\barB_X(\cA)\longrightarrow \cA
\end{equation}
is inversion on the Koszul locus. The Koszul dual algebra is obtained
by Verdier duality from the Koszul-dual coalgebra~$\cA^i$, not by
applying~$\Omega_X$ to the raw bar coalgebra.
\end{proposition}

\begin{proof}
The bar functor inserts one desuspension:
$\barB_X(\cA)=T^c(s^{-1}\bar\cA)$. Thus a finite dual of a primitive
bar generator satisfies
\[
(s^{-1}V)^\vee \cong sV^\vee,
\]
while the cobar functor applies the opposite suspension to the raw
coalgebra before forming the free algebra. These two shifts cancel in
the counit~\eqref{eq:fourier-counit-not-dual}, so
$\Omega_X\barB_X(\cA)$ reconstructs~$\cA$ when the Koszul hypotheses
of Theorem~\ref{thm:bar-cobar-inversion-qi} hold.

The finite-type holonomicity hypothesis makes Verdier duality on the
Ran bar cohomology defined in the ordinary derived category. Under
Theorem~\ref{thm:bar-cobar-isomorphism-main}, this produces the
homotopy Koszul-dual factorization algebra, and formality identifies
it with the strict algebra~$\cA^!$. The derived chiral centre has a
third target: Hochschild cochains of~$\cA$. None of the five objects in
\eqref{eq:fourier-five-objects} is a substitute for another.
\end{proof}

\subsection*{Computation~(c): $X = E_\tau$,
 $\cA = \mathcal{H}_k$}

The same elements appear as in Computation~(b), but on the elliptic
curve $E_\tau$. The propagator $\eta_{12}$ is replaced by the
period-corrected propagator $\eta^{(1)}_{12}$ of
Proposition~\ref{prop:fourier-genus1-propagator}.

\begin{computation}[Heisenberg on \texorpdfstring{$E_\tau$}{E-tau}; \ClaimStatusConditional]
\label{comp:fourier-heisenberg-elliptic}
The bar element at tensor power~$2$ on $E_\tau$ is
\[
\xi^{(1)}_{12} = s^{-1}\alpha(z_1) \otimes s^{-1}\alpha(z_2)
 \otimes \eta^{(1)}_{12}.
\]
The fiberwise bar differential $\dfib$ extracts the same OPE
coefficient $k$ via the boundary pairing. But now
$\dfib^{\,2} \neq 0$: the Arnold
failure~\eqref{eq:fourier-arnold-failure} gives
\begin{equation}\label{eq:fourier-heisenberg-curvature}
[\dfib^{\,2}(\xi^{(1)}_{123})]
=2\pi i\,k\,[\omega_\tau].
\end{equation}
The curvature is the coupling of the second-order OPE coefficient to
the Arakelov class.

The raw fiberwise operator is therefore a CDG differential. Let
$\delta_{\mathrm{per}}^{(1)}$ denote the period and boundary
correction supplied by Convention~\ref{conv:higher-genus-differentials}
and Theorem~\ref{thm:quantum-diff-squares-zero}. The total
differential is
\begin{equation}\label{eq:fourier-total-diff-genus1}
\Dg{1}=\dfib+\delta_{\mathrm{per}}^{(1)}
\end{equation}
and its defining Maurer--Cartan equation is
\begin{equation}\label{eq:fourier-period-mc}
\dfib^{\,2}+[\dfib,\delta_{\mathrm{per}}^{(1)}]
+(\delta_{\mathrm{per}}^{(1)})^2=0.
\end{equation}
\end{computation}

\begin{proposition}[\ClaimStatusConditional]
\label{prop:fourier-total-diff-nilpotent}
$\Dg{1}^{\,2} = 0$.
\end{proposition}

\begin{proof}
Expanding \eqref{eq:fourier-total-diff-genus1} gives exactly
\eqref{eq:fourier-period-mc}. The obstruction term is the class
in~\eqref{eq:fourier-heisenberg-curvature}; the period correction is
constructed so that its commutator and quadratic part represent the
negative of that class. This is the genus-$1$ specialization of
Theorem~\ref{thm:quantum-diff-squares-zero}. A single
$A$-cycle integral is not the theorem: the square-zero statement
depends on the full period-corrected modular-operadic differential.
\end{proof}

\begin{proposition}[Poincar\'e connection seed; \ClaimStatusProvedHere]
\label{prop:fourier-mukai-identification}
\index{Fourier--Mukai transform}
Let $\mathcal{P}$ be the normalized Poincar\'e bundle on
$\operatorname{Jac}(E_\tau) \times
\widehat{\operatorname{Jac}}(E_\tau)$, written in the theta
trivialization
\begin{equation}\label{eq:fourier-mukai-identification}
\mathcal{P}_{(z,\hat w)}
=\frac{\theta_1(z-\hat w|\tau)\theta'_1(0|\tau)}
{\theta_1(z|\tau)\theta_1(\hat w|\tau)} .
\end{equation}
On the complement of the theta divisors,
\begin{equation}\label{eq:fourier-poincare-connection}
d_z\log\mathcal{P}_{(z,\hat w)}
=d_z\log\theta_1(z-\hat w|\tau)
-d_z\log\theta_1(z|\tau).
\end{equation}
Thus the rank-one Heisenberg bar propagator is the infinitesimal
connection seed of the Poincar\'e kernel. The raw factorization
coalgebra $\barB_{E_\tau}(\mathcal{H}_1)$ is not a line bundle on
$\operatorname{Jac}(E_\tau)\times\widehat{\operatorname{Jac}}(E_\tau)$;
the Fourier--Mukai kernel appears only after finite-type abelianization,
Verdier duality, exponentiation of the connection, and normalization.
\end{proposition}

\begin{proof}
Equation~\eqref{eq:fourier-poincare-connection} is obtained by
logarithmically differentiating
\eqref{eq:fourier-mukai-identification} in the $z$-variable; the
factors $\theta'_1(0|\tau)$ and $\theta_1(\hat w|\tau)$ are constant
in~$z$. The first term is the same theta propagator that enters
\eqref{eq:fourier-genus1-propagator}, translated by~$\hat w$; the
second term is the normalization forcing triviality along
$\{0\}\times\widehat{E}_\tau$.
Mukai's Fourier--Mukai transform uses this line bundle as a kernel
(Mukai~\cite{Muk81}; Polishchuk~\cite{Pol03}). The Heisenberg vertex
algebra supplies the Fock realization of the same abelian current
data (Frenkel--Ben-Zvi~\cite{FBZ04}), but the category and object are
different: a Ran-space bar coalgebra becomes a Poincar\'e line bundle
only after the geometric operations stated above.
\end{proof}

\subsection*{Degeneration of the Fourier seed}
\label{subsec:fourier-degeneration}
\index{Fourier transform!recovery from bar-cobar}
\index{Fourier--Mukai transform!degeneration to Fourier}
\index{Poincar\'e line bundle!degeneration}

\begin{proposition}[Degeneration of the propagator;
 \ClaimStatusProvedHere]
\label{prop:fourier-propagator-degeneration}
Write $q = e^{2\pi i\tau}$.
The Jacobi triple product
\[
\theta_1(z|\tau)
= 2q^{1/8}\sin(\pi z)
 \prod_{n=1}^{\infty}
 (1-q^n)(1-q^n e^{2\pi iz})(1-q^n e^{-2\pi iz})
\]
gives the propagator on $E_\tau$ by logarithmic differentiation:
\begin{equation}\label{eq:fourier-theta-product}
d_z\log\theta_1(z|\tau)
= \pi\cot(\pi z)\,dz
 + 4\pi\sum_{n=1}^{\infty}
 \frac{q^n}{1-q^n}\sin(2\pi n z)\,dz.
\end{equation}
Three successive degenerations:
\begin{enumerate}
\item \emph{$q \to 0$ \textup{(}$B$-cycle collapse\textup{)}.}
 Every factor $(1 - q^n\,\cdot) \to 1$, and
 \begin{equation}\label{eq:fourier-propagator-degen}
 d_z\log\theta_1(z|\tau)
 \;\xrightarrow{\;q \to 0\;}\;
 \pi\cot(\pi z)\,dz
 = d\log\sin(\pi z).
 \end{equation}
\item \emph{Mittag-Leffler \textup{(}real locus of $\mathbb{C}^*$\textup{)}.}
 On $x \in \mathbb{R}/\mathbb{Z}$,
 \begin{equation}\label{eq:fourier-propagator-real}
 \pi\cot(\pi x)\,dx
 = \operatorname{p.v.}\!\sum_{n=-\infty}^{\infty}
 \frac{dx}{x - n}.
 \end{equation}
\item \emph{Universal cover \textup{(}$\mathbb{R} \to \mathbb{R}/\mathbb{Z}$\textup{)}.}
 Drop the $n \neq 0$ translates:
 \begin{equation}\label{eq:fourier-propagator-R}
 \frac{dx}{x} = d\log x.
 \end{equation}
\end{enumerate}
\end{proposition}

\begin{proof}
Differentiate $\log\theta_1$
term by term using $\frac{d}{dz}\log(1-q^ne^{2\pi iz})
= -\frac{2\pi i\,q^n e^{2\pi iz}}{1-q^n e^{2\pi iz}}$
and combine geometric series
(Whittaker--Watson~\cite[\S21.3]{WW27}).
The Mittag-Leffler expansion
$\pi\cot(\pi x)
= \lim_{N\to\infty}\sum_{n=-N}^{N}(x-n)^{-1}$
follows from the Weierstrass product
$\sin(\pi x) = \pi x\prod_{n=1}^{\infty}(1-x^2/n^2)$
by logarithmic differentiation
(Ahlfors~\cite[\S5.2]{Ahlfors79}).
\end{proof}

\begin{proposition}[Degeneration of the Poincar\'e connection;
 \ClaimStatusProvedHere]
\label{prop:fourier-poincare-degeneration}
\index{Poincar\'e line bundle!degeneration to character}
Let $\mathcal{P}$ be as in
\eqref{eq:fourier-mukai-identification}. Its logarithmic connection
form in the first variable is
\[
\nabla_z\log\mathcal P
=d_z\log\theta_1(z-\hat w|\tau)-d_z\log\theta_1(z|\tau).
\]
It has three controlled degenerations:
\begin{enumerate}
\item \emph{$q \to 0$.}
 The Jacobi triple product gives
 \begin{equation}\label{eq:fourier-poincare-degen-Gm}
 \nabla_z\log\mathcal P
 \;\xrightarrow{\;q \to 0\;}\;
 \pi\bigl(\cot\pi(z-\hat w)-\cot\pi z\bigr)\,dz .
 \end{equation}

\item \emph{Real torus.}
 Set $z = x$, $\hat{w} = \xi$ with
 $x,\xi \in \mathbb{R}/\mathbb{Z}$. The Mittag--Leffler expansion
 gives the principal-value distribution
 \begin{equation}\label{eq:fourier-sine-product}
 \pi\bigl(\cot\pi(x-\xi)-\cot\pi x\bigr)\,dx
 = \operatorname{p.v.}\!\sum_{n\in\mathbb Z}
 \left(\frac{1}{x-\xi-n}-\frac{1}{x-n}\right)dx .
 \end{equation}
 Exponentiating the flat integral periods along the integral
 character lattice gives the characters $e^{2\pi i n x}$,
 $n\in\mathbb Z$.

\item \emph{Universal cover.}
 $\widehat{\mathbb{R}/\mathbb{Z}} = \mathbb{Z}$;
 $\widehat{\mathbb{R}} = \mathbb{R}$. The characters
$e^{2\pi inx}$ for $n \in \mathbb{Z}$ extend to
$e^{2\pi i\xi x}$ for $\xi \in \mathbb{R}$
 \textup{(}Pontryagin~\cite{Pont66}\textup{)}, and
 \begin{equation}\label{eq:fourier-integral-from-FM}
 \hat{f}(\xi)
 = \int_{\mathbb{R}} f(x)\,e^{-2\pi i\xi x}\,dx.
 \end{equation}
 This last line uses Haar measure, an $L^1/L^2$ completion, and the
 oscillatory normalization $x\xi\mapsto -2\pi i x\xi$; these are
 analytic choices, not part of the raw bar coalgebra.
\end{enumerate}
\end{proposition}

\begin{proof}
The $q \to 0$ limit applies the Jacobi triple product to each
theta factor in~\eqref{eq:fourier-mukai-identification} and then
logarithmically differentiates in~$z$
(Polishchuk~\cite[\S16.2]{Pol03}). The distribution identity
\[
\pi\cot\pi u=\operatorname{p.v.}\!\sum_{n\in\mathbb Z}\frac{1}{u-n}
\]
is the logarithmic derivative of the Weierstrass product for
$\sin\pi u$ (Ahlfors~\cite[\S5.2]{Ahlfors79}). Pontryagin duality for
locally compact abelian groups (Rudin~\cite[\S1.7]{Rudin62}) gives
the passage from integral characters to the real-line Fourier kernel
after the analytic completions stated in item~(3).
\end{proof}

\begin{theorem}[Algebraic seed of the Fourier kernel;
 \ClaimStatusProvedHere]
\label{thm:fourier-recovery}
\index{Fourier transform!recovery from chiral bar}
For the rank-one Heisenberg Koszul contraction, the primitive
bar class $\xi=s^{-1}\alpha^*$ and the cobar generator~$x$ satisfy
\begin{equation}\label{eq:pairing-symmetric-powers}
\bigl\langle \xi^n,\,x^m\bigr\rangle=n!\,\delta_{nm}.
\end{equation}
Hence the exponential generating series of the bar-cobar pairing is
\begin{equation}\label{eq:pairing-exponential}
\sum_{n=0}^{\infty}
 \frac{\langle\xi^n,x^n\rangle}{(n!)^2}\,t^n=e^t.
\end{equation}
After the external analytic normalization $t=-2\pi i x\xi$, this
formal kernel becomes the usual Fourier kernel
$e^{-2\pi i x\xi}$. The analytic Fourier transform on
$L^2(\mathbb R)$ also requires Haar measure and topological
completion. The bar construction supplies the algebraic exponential
seed, not those analytic choices.
\end{theorem}

\begin{proof}
By Theorem~\ref{thm:fourier-heisenberg-bar}, the Koszul contraction of
$\barB(\mathcal{H}_1)$ is the symmetric coalgebra on one primitive
$\xi=s^{-1}\alpha^*$. Let $x$ be the corresponding cobar generator.
The comultiplication
$\Delta(\xi^n) = \sum_{p+q=n}\binom{n}{p}\xi^p \otimes \xi^q$
is the binomial coproduct. Multiplicativity of the bar-cobar pairing
therefore gives \eqref{eq:pairing-symmetric-powers}, and the
exponential series \eqref{eq:pairing-exponential} follows by summing
\(\sum t^n/n!\). Proposition~\ref{prop:fourier-poincare-degeneration}
identifies how the elliptic Poincar\'e connection degenerates to the
character kernel after real-form and Pontryagin-duality choices. The
present theorem is only the algebraic coalgebra calculation.
\end{proof}

\begin{remark}[Connection and holonomy]\label{rem:propagator-vs-kernel}
The propagator $d\log(z-w)$ is the Arnold/bar one-form on $C_2(X)$.
The character kernel is holonomy of the Poincar\'e connection after a
real form and a dual character have been chosen. The coalgebra
structure of $\barB(\mathcal{H}_1)$ accounts for the exponential
series in the primitive pairing; it does not by itself choose a Haar
measure, an $L^2$ completion, or the oscillatory factor $-2\pi i$.
\end{remark}

\begin{remark}[Four related kernels]\label{rem:four-transforms}
\index{Fourier transform!four incarnations}
\begin{center}
\small
\renewcommand{\arraystretch}{1.3}
\begin{tabular}{llll}
\textbf{Construction}
 & \textbf{Domain}
 & \textbf{Kernel}
 & \textbf{Extra structure} \\
\hline
Chiral bar seed
 & $\cA \mapsto \barB_X(\cA)$
 & $d\log E(z,w)$
 & chiral product on $\Ran(X)$ \\
Fourier--Mukai
 & $D^b(A) \to D^b(\hat{A})$
 & $\mathcal{P}$
 & abelian variety, Poincar\'e bundle \\
Fourier series
 & $L^2(S^1) \to \ell^2(\mathbb{Z})$
 & $e^{2\pi inx}$
 & Haar measure, integral lattice \\
Fourier integral
 & $L^2(\mathbb{R}) \to L^2(\mathbb{R})$
 & $e^{-2\pi i\xi x}$
 & Haar measure, real Pontryagin dual
\end{tabular}
\end{center}
The arrows between the rows are degeneration plus added structure.
They are not equivalences of categories, and the analytic Fourier
inversion theorem is not a formal consequence of bar-cobar inversion.
\end{remark}

\begin{proposition}[Factorization homology scope; \ClaimStatusConditional]
\label{prop:fourier-factorization-homology-scope}
Let $\cA$ be locally constant or constructible in the sense required
by the comparison theorem for factorization homology, and assume the
finite-type compactness hypotheses under which Verdier duality and
Ran chains commute with the bar construction. Then the Ran assembly of
the geometric bar coalgebra gives a comparison quasi-isomorphism
\begin{equation}\label{eq:fourier-factorization-homology}
C_\bullet^{\Ran}\bigl(X,\barB_X(\cA)\bigr)
\xrightarrow{\;\sim\;}
\int_X \cA
\end{equation}
as in Proposition~\ref{prop:bar-fh} and
Ayala--Francis~\cite{AF15}. Without these hypotheses,
$\barB_X(\cA)$ remains a factorization coalgebra on $\Ran(X)$; it is
not automatically the factorization homology of~$\cA$.
\end{proposition}

\begin{proof}
The bar construction is defined before integration: it is assembled
from logarithmic forms on $\overline{C}_n(X)$ and the chiral product.
Factorization homology is the colimit of the factorization algebra
over disks. The comparison map
\eqref{eq:fourier-factorization-homology} is the de~Rham/Ran-chain
comparison under the constructibility and compactness hypotheses of
Proposition~\ref{prop:bar-fh}. If those hypotheses fail, the colimit
defining $\int_X\cA$ need not commute with the completions and Verdier
duals used by the bar complex, so no unconditional identification is
available.
\end{proof}

\subsection*{Computation~(d): $X = \mathbb{P}^1$,
 $\cA = \widehat{\mathfrak{g}}_k$}

The affine Kac--Moody algebra $\widehat{\mathfrak{g}}_k$ at
level $k$ has currents $J^a(z)$ ($a = 1, \ldots, \dim\mathfrak{g}$)
with OPE
\begin{equation}\label{eq:fourier-km-ope}
J^a(z)J^b(w) \sim \frac{k\kappa^{ab}}{(z-w)^2}
+ \frac{f^{ab}{}_c\,J^c(w)}{z-w}
\end{equation}
where $\kappa^{ab}$ is the Killing form and $f^{ab}{}_c$ are the
structure constants of $\mathfrak{g}$.

\begin{computation}[Kac--Moody bar; \ClaimStatusProvedHere]
\label{comp:fourier-km-bar}
At tensor power~$2$, the bar element
$s^{-1}J^a(z_1) \otimes s^{-1}J^b(z_2) \otimes \eta_{12}$ receives
two contributions from the bar differential:
\begin{enumerate}
\item \emph{First-order pole.}
 $\bigl\langle f^{ab}{}_c\, J^c(w)/(z-w),\;
 \eta_{12}\bigr\rangle_{D_{12}}
 = f^{ab}{}_c \cdot s^{-1}J^c$:
 the Lie bracket term.
\item \emph{Second-order pole.}
 $\bigl\langle k\kappa^{ab}/(z-w)^2,\;
 \eta_{12}\bigr\rangle_{D_{12}}
 = k\kappa^{ab} \cdot 1$:
 the Killing term.
\end{enumerate}
The first contribution builds the Chevalley--Eilenberg differential;
the second introduces curvature.
\end{computation}

\begin{theorem}[\ClaimStatusProvedHere]
\label{thm:fourier-km-bar}
\index{Kac--Moody!bar complex}
The bar complex of $\widehat{\mathfrak{g}}_k$ on $\mathbb{P}^1$ is
the curved chiral Chevalley--Eilenberg coalgebra:
\begin{equation}\label{eq:fourier-km-bar}
\barB(\widehat{\mathfrak{g}}_k) \;\cong\;
\hat{C}^\bullet_{\mathrm{ch}}(\mathfrak{g},\, k\kappa).
\end{equation}
The curvature element $m_0 = k\kappa \in \barB_0$ is the Killing
form at level~$k$. The affine level $-k-2h^\vee$ belongs to the
post-Verdier Koszul-dual algebra on the finite-type noncritical lane;
it is not a parameter of the raw bar coalgebra.
\end{theorem}

\begin{proof}
The first-order pole contributes the Lie coalgebra differential
$d_{\mathrm{CE}}\colon s^{-1}\mathfrak{g} \otimes
s^{-1}\mathfrak{g} \to s^{-1}\mathfrak{g}$, giving the
Chevalley--Eilenberg complex $C^\bullet(\mathfrak{g})$. The
second-order pole contributes the curvature $m_0 = k\kappa$, making
the complex curved: $m_1^2(x) = m_2(m_0, x) - m_2(x, m_0) =
[m_0, x]$ (the commutator, with minus sign).

After taking bar cohomology and applying Verdier finite duality, the
dual affine parameter on the noncritical finite-type lane is
$k \mapsto -k - 2h^\vee$. Thus
$\kappa(\widehat{\mathfrak{g}}_k^!)
= \kappa(\widehat{\mathfrak{g}}_{-k-2h^\vee})$ in the vertex-algebra
normalization. The Koszul dual
$(\widehat{\mathfrak{g}}_k)^!
= \mathrm{CE}^{\mathrm{ch}}(\widehat{\mathfrak{g}}_{-k-2h^\vee})$
is the chiral CE algebra at the dual level, not the raw affine
algebra $\widehat{\mathfrak{g}}_{-k-2h^\vee}$ itself. The
Feigin--Frenkel involution~\cite{Feigin-Frenkel} operates within the
affine family; Koszul duality changes the categorical target.
See \S\ref{sec:km-koszul-abstract}.
\end{proof}


% ================================================================
\section{Specialization to a point}\label{sec:fourier-specialization}
% ================================================================

\begin{theorem}[Specialization; \ClaimStatusProvedHere]
\label{thm:fourier-specialization}
\index{bar construction!specialization to a point}
Let $X = \operatorname{Spec} k$. Then:
\begin{enumerate}
\item $\overline{C}_n(\operatorname{Spec} k) = \operatorname{Spec} k$
 for all $n$.
\item $\Omega^\bullet(\overline{C}_n,\, \log D) = k$ (no forms, no
 propagator).
\item $d_{\mathrm{dR}} = 0$ and $d_{\mathrm{res}}$ reduces to the
 binary multiplication, bracket, or operadic structure map appropriate
 to the algebraic input.
\item $\barB_X(\cA)$ reduces to
 $T^c(s^{-1}\bar{\cA},\, d_{\mathrm{bracket}})$, the classical bar
 construction of Loday--Vallette~\cite[\S2.2]{LV12}.
\end{enumerate}
\end{theorem}

\begin{proof}
Items~(1)--(3) are immediate: $\operatorname{Spec} k$ has no
nontrivial configuration geometry. For~(4), the general
three-component differential
$d = d_{\mathrm{int}} + d_{\mathrm{res}} + d_{\mathrm{dR}}$
collapses to $d_{\mathrm{int}} + d_{\mathrm{res}}$ when
$X = \operatorname{Spec} k$. The residue component
$d_{\mathrm{res}}$ at $X = \operatorname{Spec} k$ has no
configuration-space residue left; it is the algebraic binary
operation of the chosen operadic ambient. The result is the classical
bar construction
$T^c(s^{-1}\bar{\cA},\, d_2)$ with the standard bar differential.
\end{proof}

\begin{equation}\label{eq:fourier-specialization-diagram}
\begin{tikzcd}[column sep=huge, row sep=large]
\mathrm{ChirAlg}(X) \arrow[r, "\barB_X"]
 \arrow[d, "X = \mathrm{pt}"']
& \mathrm{Fact}^c(\Ran\, X) \arrow[d, "X = \mathrm{pt}"] \\
\mathrm{Alg}_k \arrow[r, "\barB"']
& \mathrm{CoAlg}_k
\end{tikzcd}
\end{equation}

\begin{remark}\label{rem:fourier-zero-dim-shadow}
Classical operadic Koszul duality
(Loday--Vallette~\cite{LV12}) is the zero-dimensional specialization
of the chiral bar-cobar construction. A theorem lifts from the point
to curves only after the relevant conilpotence, finite-type,
holonomicity, and logarithmic compactification hypotheses have been
installed; the configuration-space geometry supplies the new
obstruction classes rather than a formal automatic lift.
\end{remark}


% ================================================================
\section{The four properties}\label{sec:fourier-four-properties}
% ================================================================

\begin{theorem}[Four precise Fourier analogues;
\ClaimStatusConditional]
\label{thm:fourier-four-properties}
The Fourier language in this chapter has the following precise
mathematical content:
\begin{enumerate}[label=\textup{(\Alph*)}]
\item \emph{Verdier seed.}
 On the finite-type Verdier/Koszul lane,
 \[
 \mathbb{D}_{\Ran}\barB_X(\cA)\simeq \cA^!_\infty
 \]
 \textup{(}Theorem~\ref{thm:bar-cobar-isomorphism-main}\textup{)}.
 If the homotopy dual is formal, this identifies with the strict
 algebra~$\cA^!$. The target is the post-Verdier dual algebra, not
 the raw bar coalgebra $\barB_X(\cA^!)$.
\item \emph{Inversion.}
 On the strict conilpotent Koszul locus,
 \[
 \Omega_X\barB_X(\cA)\xrightarrow{\;\sim\;}\cA
 \]
 \textup{(}Theorems~\ref{thm:bar-cobar-inversion-qi},
 \ref{thm:higher-genus-inversion}\textup{)}. Off that locus the
 strongest proved statement is completed or coderived, not ordinary
 analytic Fourier inversion.
\item \emph{Poincar\'e--Parseval shadow.}
 Under the BV-perfect, Verdier-nondegenerate, uniform-weight-perfect
 hypotheses of Theorem~\ref{thm:quantum-complementarity-main},
 $Q_g(\cA) \oplus Q_g(\cA^!) \simeq
 R\Gamma(\overline{\mathcal{M}}_g,\, Z_{\cA})$ for $g\ge1$, with the
 genus-$0$ case treated separately. This is the Verdier/Lagrangian
 shadow of Parseval, not a norm identity on a Hilbert space.
\item \emph{Characteristic function.}
 On the uniform-weight scalar-diagonal lane,
 \[
 \mathrm{obs}_g(\cA)=\kappa(\cA)\lambda_g.
 \]
 At the all-weight numerical level,
 $F_g(\cA)=\kappa(\cA)\lambda_g^{\mathrm{FP}}
 +\delta F_g^{\mathrm{cross}}(\cA)$, with
 $\delta F_1^{\mathrm{cross}}=0$
 \textup{(}Theorem~\ref{thm:modular-characteristic}\textup{)}. The
 $\hat A$-series is the scalar-lane generating function.
\end{enumerate}
\end{theorem}

\begin{remark}[Classical Fourier analogues with scope]
\label{rem:fourier-classical-analogues}
\begin{center}
\small
\begin{tabular}{lll}
\textbf{Property}
 & \textbf{Chiral (this monograph)}
 & \textbf{Classical Fourier} \\
\hline
(A) & Verdier seed
 & $\widehat{f \cdot g} = \hat{f} * \hat{g}$ \\
(B) & Bar-cobar inversion
 & Fourier inversion after measure and topology \\
(C) & Verdier/Lagrangian complementarity
 & $\|f\|^2 = \|\hat{f}\|^2$ (Parseval) \\
(D) & Modular characteristic $\kappa$
 & Characteristic function determines moments
\end{tabular}
\end{center}
Only the chiral column is proved inside the bar-cobar theory. The
classical column is a dictionary of analogues after analytic structures have
been added.
\end{remark}

\begin{remark}\label{rem:fourier-scope}
These four properties, the specialization theorem
(\S\ref{sec:fourier-specialization}), and the chiral Hochschild polynomial
growth theorem (Theorem~H;
Theorem~\ref{thm:hochschild-polynomial-growth})
constitute the structural core.
Part~\ref{part:bar-complex} develops the proofs. Part~\ref{part:physics-bridges} computes the transform for each
standard family.
\end{remark}


% ================================================================
\section{The landscape}\label{sec:fourier-landscape}
% ================================================================

\begin{center}
\small
\begin{tabular}{llll}
\textbf{$\cA$}
 & \textbf{$\barB(\cA)$}
 & \textbf{$\kappa(\cA)$}
 & \textbf{$c(\cA)$} \\
\hline
$\mathcal{H}_k$
 & $\operatorname{Sym}^{\mathrm{ch},c}(V^*[-1])$; dual curvature $-k$
 & $k$
 & $1$ \\[2pt]
$\widehat{\mathfrak{g}}_k$
 & Curved CE coalgebra; post-Verdier level $-k - 2h^\vee$
 & $\dfrac{(k + h^\vee)\dim\mathfrak{g}}{2h^\vee}$
 & $\dfrac{k\dim\mathfrak{g}}{k + h^\vee}$ \\[6pt]
$\mathrm{Vir}_c$
 & Curved CE coalgebra
 & $c/2$
 & $c$ \\[2pt]
$\mathcal{W}_N$
 & Via DS reduction
 & DS formula
 & DS formula
\end{tabular}
\end{center}

\noindent
The detailed computations appear in Part~\ref{part:physics-bridges}
(Chapters~\ref{ch:heisenberg-frame}--\ref{chap:yangians}).


% ================================================================
\section{Genus as deformation variable}\label{sec:fourier-genus}
% ================================================================

\begin{remark}[Ref.\
 Theorem~\ref{thm:mc2-full-resolution}]
\label{rem:fourier-genus-preview}
At genus $g \geq 1$, the bar complex $\barB_g(\cA)$ is a family of
curved cochain complexes over $\overline{\mathcal{M}}_g$. The proved
invariant $\kappa(\cA)$
(Theorem~\ref{thm:modular-characteristic}) is the scalar shadow of the
universal Maurer--Cartan class
\begin{equation}\label{eq:fourier-theta-preview}
\Theta_{\cA} \;\in\;
\mathrm{MC}\!\left(
 \Defcyc(\cA) \;\hat{\otimes}\;
 R\Gamma(\overline{\mathcal{M}}_{g,\bullet},\, \mathbb{Q})
\right),
\end{equation}
constructed by Theorem~\ref{thm:mc2-full-resolution}
(Theorem~\ref{thm:master-theta}).
On the theorem surface this means that the Fourier seed lands in the
concrete modular datum of
Theorem~\ref{thm:volume-one-concrete-modular-datum}: its M-level
output is the curved family $(\barB_g(\cA), \Dg{g})$, its S-level
shadow is the characteristic hierarchy, and its H-level logarithm
is~$\Theta_{\cA}$ (constructed bar-intrinsically by
Theorem~\ref{thm:mc2-bar-intrinsic}: $\Theta_{\cA} := D_{\cA} - d_0$
is automatically MC because $D_{\cA}^2 = 0$, with no restriction to
simple Lie symmetry). What remains open is the coderived
continuation beyond the proved algebraic loci, not the existence of
this datum.
\end{remark}
